% Caderno 5 — projeção da meia fração num fatorial completo com réplicas.
% Base: Peixoto, `Capítulo 4.pdf`, p. 28 a 31.
% Plano 2^{3-1}_{III} com geratriz I = 123, quatro ensaios:
%   1 (-,-,+)   2 (+,-,-)   3 (-,+,-)   4 (+,+,+)
% Com um fator inerte, os quatro ensaios projetam-se num 2^2 completo;
% com dois fatores inertes, projetam-se em dois níveis com duas réplicas
% cada. É o argumento mais forte a favor da triagem.
\documentclass[border=3pt,11pt]{standalone}
% Preâmbulo comum das figuras TikZ das apostilas.
% Carregado por todos os `standalone` deste diretório, para que os cinco
% diagramas não divirjam em fonte, cor ou espessura de traço.
%
% fontenc T1 + inputenc utf8 são obrigatórios: sem eles "ç", "ã" e "é"
% saem trocados ou desaparecem silenciosamente do PDF.
\usepackage[T1]{fontenc}
\usepackage[utf8]{inputenc}
\usepackage[brazil]{babel}
% A ordem importa: `amssymb` antes de `newtxmath`, exatamente como em
% `doeapostila.cls`. Invertida, as duas definem `\Bbbk` e a compilação aborta.
\usepackage{amsmath,amssymb}
\usepackage{newtxtext,newtxmath}   % mesma família serifada da classe doeapostila
\usepackage{tikz}
\usetikzlibrary{arrows.meta,calc,decorations.pathreplacing,positioning,
  shapes.geometric,backgrounds}

% Paleta: preto, cinzas e uma única cor de apoio.
\definecolor{azulDoE}{HTML}{1B4E79}
\definecolor{cinzaDoE}{gray}{0.55}
\definecolor{cinzaClaro}{gray}{0.82}
\definecolor{cinzaFundo}{gray}{0.93}

\tikzset{
  % vértice de plano fatorial, com a resposta escrita dentro
  vertice/.style={circle, draw=black, fill=white, line width=0.5pt,
                  inner sep=0pt, minimum size=8.5mm, font=\scriptsize},
  vertdestaque/.style={circle, draw=azulDoE, fill=azulDoE!10, text=azulDoE,
                  line width=0.9pt, inner sep=0pt, minimum size=8.5mm,
                  font=\scriptsize},
  % ponto de planejamento, sem texto dentro
  ponto/.style={circle, draw=black, fill=black, inner sep=0pt,
                minimum size=2.4mm},
  pontovazio/.style={circle, draw=black, fill=white, line width=0.5pt,
                inner sep=0pt, minimum size=2.6mm},
  pontoazul/.style={circle, draw=azulDoE, fill=azulDoE, inner sep=0pt,
                minimum size=2.8mm},
  axial/.style={rectangle, draw=azulDoE, fill=azulDoE, inner sep=0pt,
                minimum size=2.6mm},
  apagado/.style={circle, draw=cinzaClaro, fill=white, line width=0.5pt,
                inner sep=0pt, minimum size=2.6mm},
  aresta/.style={draw=black, line width=0.5pt},
  arestafundo/.style={draw=cinzaClaro, line width=0.5pt},
  eixo/.style={draw=black, line width=0.5pt, -{Latex[length=1.8mm,width=1.4mm]}},
  seta/.style={draw=black, line width=0.6pt, -{Latex[length=2.2mm,width=1.7mm]}},
  setaazul/.style={draw=azulDoE, line width=0.8pt,
                   -{Latex[length=2.2mm,width=1.7mm]}},
  cota/.style={draw=cinzaDoE, line width=0.5pt,
               {Latex[length=1.6mm]}-{Latex[length=1.6mm]}},
  % Corpos escolhidos para casar com as figuras de matplotlib, que usam
  % 8 pt nos ticks e 9 pt nos rótulos de eixo. Com a base de 11 pt do
  % `standalone` (a mesma da classe doeapostila), \scriptsize = 8 pt e
  % \footnotesize = 9 pt. Abaixo disso a figura deixa de ser legível na
  % página impressa, e é para isso que a folha de prova serve.
  rotulo/.style={font=\footnotesize},
  rotulinho/.style={font=\scriptsize},
  nota/.style={font=\footnotesize, align=left},
  notaazul/.style={font=\footnotesize, align=left, text=azulDoE},
}

\begin{document}
\begin{tikzpicture}[x=1cm, y=1cm]

  % Rótulo de ensaio afastado do marcador: encostado, o número some sob o disco.
  \tikzset{marca/.style={rotulinho, text=azulDoE, inner sep=0pt}}

  % ================= estágio 1: o cubo com os quatro ensaios =============
  \begin{scope}
    \def\lado{2.3}
    \def\dx{0.95}
    \def\dy{0.80}

    % f1 -> direita, f3 -> cima, f2 -> fundo
    \coordinate (c1) at (0,\lado);                 % ensaio 1: (-,-,+)
    \coordinate (c2) at (\lado,0);                 % ensaio 2: (+,-,-)
    \coordinate (c3) at (\dx,\dy);                 % ensaio 3: (-,+,-)
    \coordinate (c4) at (\lado+\dx,\lado+\dy);     % ensaio 4: (+,+,+)
    \coordinate (v1) at (0,0);
    \coordinate (v2) at (\lado,\lado);
    \coordinate (v3) at (\dx,\lado+\dy);
    \coordinate (v4) at (\lado+\dx,\dy);

    \draw[arestafundo] (c3) -- (v3) -- (c4) -- (v4) -- cycle;
    \draw[arestafundo] (v1) -- (c3);
    \draw[aresta] (v1) -- (c2) -- (v2) -- (c1) -- cycle;
    \draw[aresta] (c2) -- (v4);
    \draw[aresta] (c1) -- (v3);
    \draw[aresta] (v2) -- (c4);

    \node[apagado] at (v1) {}; \node[apagado] at (v2) {};
    \node[apagado] at (v3) {}; \node[apagado] at (v4) {};

    \node[pontoazul] at (c1) {};  \node[marca] at ($(c1)+(-0.26,0.22)$) {1};
    \node[pontoazul] at (c2) {};  \node[marca] at ($(c2)+(0.26,-0.22)$) {2};
    \node[pontoazul] at (c3) {};  \node[marca] at ($(c3)+(-0.26,-0.22)$) {3};
    \node[pontoazul] at (c4) {};  \node[marca] at ($(c4)+(0.26,0.22)$) {4};

    \coordinate (o) at (-1.15,-1.05);
    \draw[eixo] (o) -- ++(0.80,0)    node[rotulinho, right] {1};
    \draw[eixo] (o) -- ++(0,0.80)    node[rotulinho, above] {3};
    \draw[eixo] (o) -- ++(0.55,0.46) node[rotulinho, above right, xshift=-1mm] {2};

    \node[nota, anchor=north, align=center] at (1.6,-1.55)
      {$2^{3-1}_{\mathrm{III}}$, $I = 123$\\ 4 ensaios num cubo de 8 vértices};
  \end{scope}

  % ================= seta 1 ==============================================
  \begin{scope}[xshift=4.35cm]
    \draw[seta] (0,1.55) -- (1.35,1.55);
    \node[rotulinho, anchor=south, align=center] at (0.68,1.66)
      {fator 1\\ inerte};
  \end{scope}

  % ================= estágio 2: o 2^2 completo ===========================
  \begin{scope}[xshift=6.25cm, yshift=0.28cm]
    \def\l{2.0}
    \coordinate (q1) at (0,\l);      % ensaio 1: f2 -, f3 +
    \coordinate (q2) at (0,0);       % ensaio 2: f2 -, f3 -
    \coordinate (q3) at (\l,0);      % ensaio 3: f2 +, f3 -
    \coordinate (q4) at (\l,\l);     % ensaio 4: f2 +, f3 +

    \draw[aresta] (q2) -- (q3) -- (q4) -- (q1) -- cycle;
    \node[pontoazul] at (q1) {};  \node[marca] at ($(q1)+(-0.26,0.22)$) {1};
    \node[pontoazul] at (q2) {};  \node[marca] at ($(q2)+(-0.26,-0.22)$) {2};
    \node[pontoazul] at (q3) {};  \node[marca] at ($(q3)+(0.26,-0.22)$) {3};
    \node[pontoazul] at (q4) {};  \node[marca] at ($(q4)+(0.26,0.22)$) {4};

    \coordinate (o2) at (-0.80,-0.78);
    \draw[eixo] (o2) -- ++(0.70,0) node[rotulinho, right] {2};
    \draw[eixo] (o2) -- ++(0,0.70) node[rotulinho, above] {3};

    \node[nota, anchor=north, align=center] at (1.0,-1.83)
      {$2^{2}$ completo\\ um ensaio por vértice};
  \end{scope}

  % ================= seta 2 ==============================================
  \begin{scope}[xshift=9.15cm]
    \draw[seta] (0,1.55) -- (1.35,1.55);
    \node[rotulinho, anchor=south, align=center] at (0.68,1.66)
      {fatores 1 e 3\\ inertes};
  \end{scope}

  % ================= estágio 3: dois níveis com duas réplicas ============
  \begin{scope}[xshift=11.05cm, yshift=0.28cm]
    \draw[eixo] (-0.30,0) -- (2.45,0) node[rotulinho, right] {2};
    \draw[cinzaDoE, line width=0.4pt] (0.30,0.07) -- (0.30,-0.07);
    \draw[cinzaDoE, line width=0.4pt] (1.80,0.07) -- (1.80,-0.07);
    \node[rotulinho, below] at (0.30,-0.12) {$-1$};
    \node[rotulinho, below] at (1.80,-0.12) {$+1$};

    \draw[cinzaClaro, line width=0.5pt] (0.30,0.70) -- (0.30,1.50);
    \draw[cinzaClaro, line width=0.5pt] (1.80,0.70) -- (1.80,1.50);

    \node[pontoazul] at (0.30,0.70) {}; \node[marca] at (0.58,0.70) {1};
    \node[pontoazul] at (0.30,1.50) {}; \node[marca] at (0.58,1.50) {2};
    \node[pontoazul] at (1.80,0.70) {}; \node[marca] at (2.08,0.70) {3};
    \node[pontoazul] at (1.80,1.50) {}; \node[marca] at (2.08,1.50) {4};

    \node[nota, anchor=north, align=center] at (1.05,-1.55)
      {dois níveis\\ com duas réplicas cada};
  \end{scope}

  \node[nota, anchor=north west, text width=13.9cm] at (-1.15,-2.95)
    {Um fator que não afeta a resposta não some do planejamento: ele devolve
     ensaios. A fração deixa de ser fração e volta a ser fatorial completo
     nos fatores que restaram, com réplicas de brinde.};

\end{tikzpicture}
\end{document}
